\section{Discussion}
\label{sec:discussion}

\subsection{Risks and Limitations}

\textbf{Smart-contract risk.} The hub, issuers, CDPs, and debt tokens share common contract logic. A critical vulnerability in a widely-used contract could affect multiple markets. Formal verification covers token conservation and epoch monotonicity; broader functional coverage is an ongoing effort.

\textbf{Oracle staleness.} Liquidation detection depends on timely oracle updates. If oracle prices lag during a flash crash, detection is delayed and recovery proceeds may fall below debt. The two-epoch structure bounds this lag to one additional epoch window; markets with more volatile collateral should configure shorter epochs.

\textbf{Proof-system failures.} If the ZK prover becomes unavailable, no new epochs can advance. Existing positions are safe (funds remain in CDPs) but cannot settle new interest or originate. Recovery requires a new prover deployment. No fallback to optimistic settlement is currently implemented.

\textbf{Dual-liquidation race.} A liquidity-forwarding loan exposes the vault to the gap between \daloc{}'s two-epoch liquidation and the venue's instant liquidation trigger. Reserve sizing (Section~\ref{sec:risk}) attempts to make this gap negligible, but a sufficiently sharp collateral crash can trigger the venue's liquidation before \daloc{}'s grace period completes, leaving a residual shortfall absorbed by the vault.

\textbf{Cross-chain liveness.} Bridge delays or outages stall capital delivery on Prime Trade origination and close. Positions remain valid but economically frozen until the bridge resumes. Markets that rely on a single bridge path inherit its liveness profile.

\textbf{Host-chain downtime.} \daloc{} markets inherit the liveness of the chains on which their contracts are deployed. A chain halt delays epoch advances, hook execution, and liquidations while economic risk continues to move elsewhere. Markets on L2s additionally inherit the rollup's recovery properties. Chain downtime is a modeled liveness risk (Section~\ref{sec:settlement}), not a correctness failure.

\textbf{External-venue liquidity and liveness.} Liquidity-forwarding loans depend on external venues for borrowing, repayment, and liquidation timing. If a venue pauses, has insufficient liquidity, or liquidates faster than \daloc{}'s two-epoch process, the loan may enter a pending-unwind state and residual losses flow through the loss waterfall (Section~\ref{sec:risk}). Reserve sizing makes this risk explicit and priced but cannot remove it.

\textbf{Prime Trade venue downtime.} Prime Trade positions depend on HyperEVM and HyperCore remaining live. If the destination venue stalls, the margin agent cannot rebalance, stop out, or recall capital on schedule. The position remains constrained by its program, but economic exposure changes while execution is unavailable. Venue downtime is modeled in the Risk Kernel and Risk House policy.

\textbf{Stablecoin and bridge centralization.} For USDC-based markets, the principal asset issuer (Circle) can freeze balances at specific addresses. If a CDP, vault, or margin agent address is blacklisted, the affected positions cannot settle normally. This concentration risk is mitigated by market diversity (alternative principal assets and bridge adapters may serve separate markets) but cannot be eliminated while a single issuer's token is the primary path.

\subsection{Future Capabilities}

The architecture is designed for staged extension. Near-term additions include:
\begin{itemize}
  \item \textbf{Risk House ecosystem.} The architecture already admits permissionless Risk House registration (Section~\ref{sec:risk}); at scale, multiple independent houses compete on terms and backstop specific collateral classes. The kernel bounds prevent a weak house from degrading system-wide safety.
  \item \textbf{Under-collateralized credit.} On-chain reputation, verifiable off-chain creditworthiness (e.g.\ institutional attestations), and cross-protocol track records enable LTV thresholds above~100\% for qualified borrowers under tighter Risk House oversight.
  \item \textbf{Cross-margining.} Correlated positions under the same borrower and same Risk House can share margin via a portfolio-level health factor, improving capital efficiency without weakening per-position isolation at the settlement layer.
  \item \textbf{Options and nonlinear collateral.} Collateral-class adapters for options and other nonlinear claims, valued at stressed executable close-out value rather than oracle mark, extend the protocol beyond spot and yield-bearing tokens.
  \item \textbf{Multi-venue Prime Trade.} Margin agents generalize beyond HyperEVM to additional perpetual venues, each as a separate bridge adapter and constraint set.
  \item \textbf{Operator ecosystem.} The architecture supports multiple independent vault operators, each running a vault instance with a distinct mandate (one conservative with high-grade collateral and short tenors, another aggressive with broader collateral and higher concentration limits). Competition between operators for LP capital creates pressure toward better risk-adjusted returns. Per-market settlement operators may independently rotate or compete, mirroring vault-level competition at the settlement layer.
  \item \textbf{Institutional access.} Criteria-based matching (Section~\ref{sec:architecture}) already enables regulated counterparties to participate without a permissioned fork. Realizing this in production requires only attestation issuers and criteria definitions, not protocol modification.
  \item \textbf{Maturity management.} Refinancing (Section~\ref{sec:lifecycle}) already allows mid-term rollovers. At scale, automated maturity laddering (where the operator pre-matches refinancing quotes ahead of expiry) can smooth redemption waves, reduce cliff-event risk, and improve capital utilization across the loan book.
\end{itemize}

\subsection{Conclusion}

\daloc{} unifies collateralized credit, cross-chain capital coordination, and risk quantification within a single settlement layer. The Harvest Vault funds loans directly under a proof-attested mandate; external venues provide overflow capacity beyond the vault's balance sheet. Fixed-rate, fixed-term loans; leveraged perpetual trading without selling the underlying; and passive yield from a portfolio of fixed-rate loans share a common matching and settlement infrastructure.

Criteria-based matching lets restricted and unrestricted counterparties share a single liquidity surface. Zero-knowledge proofs enforce settlement correctness and underwriting integrity; mandate compliance is enforced through fee-gated compliance proof. Per-position and per-market isolation bound the blast radius of defaults. The protocol presents an architecture for the transition from pool-based, variable-rate, single-chain lending to structured, provable, multi-chain credit.
